% =====================================================================================
% RASCUNHO 2026-09-26 — Reconhecimento do número: controlo pela confiança da leitura
% Números marcados com [P] são PROVISÓRIOS (117 vídeos); trocar pelos finais (159) amanhã.
% Onde entra cada bloco está indicado no cabeçalho de cada um.
% =====================================================================================


% -------------------------------------------------------------------------------------
% BLOCO A — Cap. 3 (Chapter3_final.tex). SUBSTITUI o 2.º parágrafo de
% "Recorte da Região Dorsal por Estimativa de Pose" (o do filtro de qualidade).
% -------------------------------------------------------------------------------------

Antes de prosseguir, cada recorte passa por um filtro de qualidade que descarta os casos sem
informação utilizável. São rejeitados os recortes com dimensões inferiores a $15\times20$ píxeis,
tipicamente jogadores muito distantes da câmara, e os recortes desfocados, quantificados pela
variância do Laplaciano (limiar de $15$). Estes limiares são deliberadamente permissivos. Numa
primeira versão do módulo, o filtro de qualidade era a única proteção contra leituras de recortes
sem informação e tinha de ser exigente ($25\times30$ píxeis e nitidez de $40$), à custa de excluir
jogadores cujo número era, ainda assim, legível. Com o controlo pela confiança da leitura descrito
na Secção~\ref{sec:met-jnr-confianca}, é o próprio modelo que indica quando não conseguiu ler, e o
filtro passa a ter apenas a função de evitar chamadas inúteis ao modelo.


% -------------------------------------------------------------------------------------
% BLOCO B — Cap. 3. SUBSTITUI a subsecção "Agregação Temporal por Votação" inteira
% (mantém a figura fig:jnr-crops, que continua a fazer sentido).
% -------------------------------------------------------------------------------------

\subsection{Confiança da Leitura e Decisão por Janela}\label{sec:met-jnr-confianca}

A leitura efetuada sobre um recorte isolado é, por natureza, ruidosa. Uma pose desfavorável, uma
oclusão momentânea ou variações de iluminação são suficientes para induzir erros na inferência de
um \emph{frame}, mesmo quando a numeração se mantém globalmente legível ao longo do
\emph{tracklet} (Figura~\ref{fig:jnr-crops}). A decisão final tem, por isso, de combinar várias
leituras, e a forma de as combinar determina o tipo de erro que o sistema comete.

Os recortes de cada \emph{tracklet} que passam o filtro de qualidade são ordenados no tempo e
divididos em janelas de seis recortes consecutivos. Cada janela é lida pelo Qwen2.5-VL a partir
dos seus cinco recortes mais nítidos, com a mesma \emph{prompt} e a mesma descodificação
determinística descritas na secção anterior. Além da resposta, regista-se a probabilidade que o
modelo atribuiu a cada \emph{token} gerado. A confiança da janela é o produto das probabilidades
dos \emph{tokens} que formam o número. Um número de dois dígitos lido com segurança tem confiança
próxima de $1$, enquanto um número adivinhado num recorte pouco legível raramente ultrapassa
$0.3$, porque a incerteza em qualquer dos dígitos baixa o produto.

A decisão segue duas regras simples:
\begin{enumerate}
    \item uma janela só vota se tiver lido um número com confiança igual ou superior a $0.5$; as
    janelas com confiança inferior e as janelas em que o modelo recusou a leitura não votam;
    \item o \emph{tracklet} recebe o número mais votado entre as janelas que votaram; se nenhuma
    votou, o \emph{tracklet} fica sem número.
\end{enumerate}

Esta regra resulta de um diagnóstico feito ao próprio módulo. Na primeira versão, a função que
chamava o modelo devolvia sempre a mesma confiança fixa, qualquer que fosse a certeza da leitura,
e a decisão dependia só da contagem de votos, com uma quota mínima de concordância entre janelas.
Um número nítido e um número adivinhado num recorte desfocado pesavam o mesmo. Quando um jogador
era visto durante muito tempo em más condições, as leituras erradas acumulavam votos suficientes
para vencer, e o erro propagava-se a centenas de deteções. A confiança do modelo já existia, mas
era descartada. Usá-la permite separar as leituras seguras das restantes antes da votação, em vez
de tentar compensar a posteriori com quotas de concordância.

O limiar de $0.5$ traduz uma escolha de projeto. No SoccerNet-GSR, um número errado não só falha a
correspondência com a anotação como a atribui a uma identidade que não existe, enquanto um número
em falta apenas deixa de a atribuir. Entre errar e não responder, o sistema foi desenhado para não
responder. O custo desta escolha, medido no Capítulo~\ref{cap:resultados}, é um ligeiro aumento dos
jogadores que ficam sem número.


% -------------------------------------------------------------------------------------
% BLOCO C — Cap. 3, Tabela tab:hiperparametros. SUBSTITUIR as linhas do filtro de qualidade e
% de "leitura e votação" por estas.
% -------------------------------------------------------------------------------------

%        \multicolumn{3}{|l|}{\emph{Reconhecimento de número de camisola (filtro de qualidade)}} \\
%        \hline
%        Dimensão mínima do recorte   & $15\times20$~px & Rejeita jogadores muito distantes \\
%        \hline
%        Nitidez mínima               & $15.0$          & Variância do Laplaciano \\
%        \hline
%        \multicolumn{3}{|l|}{\emph{Reconhecimento de número de camisola (leitura e decisão)}} \\
%        \hline
%        Modelo de leitura            & Qwen2.5-VL-3B (4 bits) & Inferência do número \\
%        \hline
%        Recortes por janela          & $6$             & Unidade de decisão temporal \\
%        \hline
%        Recortes lidos por janela    & $5$             & Os mais nítidos da janela \\
%        \hline
%        Confiança mínima para votar  & $0.5$           & Produto das probabilidades dos \emph{tokens} \\
%        \hline


% -------------------------------------------------------------------------------------
% BLOCO D — Cap. 4 (gerador/Chapter4_template.tex). NOVA subsecção, a seguir a
% "Resultados" (sec:res_jnr_resultados) e ANTES de "Teto da Abordagem".
% -------------------------------------------------------------------------------------

\subsection{Controlo pela Confiança da Leitura}\label{sec:res_jnr_confianca}

A exatidão por \emph{tracklet} da Tabela~\ref{tab:jnr} trata todos os \emph{tracklets} por igual.
O GS-HOTA, pelo contrário, conta deteções, e um \emph{tracklet} longo com o número errado pesa
tanto como centenas de \emph{tracklets} curtos. Para perceber onde estava o custo do número na
métrica oficial, as deteções de jogadores emparelhadas com a anotação foram classificadas em cinco
categorias, conforme a anotação tem ou não número e o sistema atribui ou não um número
(Tabela~\ref{tab:jnr-categorias}).

Na primeira versão do módulo, sobre os $159$ vídeos, $51.4\%$ das deteções recebiam o número
certo, $12.3\%$ recebiam um número errado e $2.5\%$ recebiam um número quando o jogador não tinha
número visível na anotação. A maior parte do erro estava em \emph{tracklets} longos: $9.2\%$ de
todas as deteções pertenciam a \emph{tracklets} com mais de $200$ \emph{frames} que terminavam com
um número errado. Reconstituindo os votos de cada janela a partir dos registos de execução, em
cerca de um terço das deteções com número errado o número certo chegou a ser lido, mas perdeu a
votação para leituras erradas mais numerosas. Nas restantes, o número certo nunca foi lido, e as
leituras erradas vinham de recortes em que o número não era legível. Em ambos os casos, o que faltava à decisão era saber
quais das leituras eram seguras.

A regra de decisão pela confiança (Secção~\ref{sec:met-jnr-confianca}) foi primeiro testada numa
amostra aleatória de $10$ vídeos de teste e validação e, confirmada a tendência, alargada por fases
a amostras aleatórias sucessivas, até cobrir % [P]
$117$ dos $159$ vídeos. % [P]
Em cada fase, a primeira versão e a versão com confiança foram comparadas sobre as mesmas deteções,
com a mesma anotação e com o GS-HOTA oficial. A Tabela~\ref{tab:jnr-categorias} resume o resultado.

\begin{table}[htbp]
    \centering
    \caption{O que acontece às deteções de jogadores, em percentagem, na primeira versão do módulo e
    com a decisão pela confiança da leitura ($117$ vídeos). % [P]
    }
    \label{tab:jnr-categorias}
    \begin{tabular}{|l|c|c|c|}
        \hline
        \textbf{Categoria} & \textbf{Primeira versão} & \textbf{Com confiança} & \textbf{$\Delta$} \\
        \hline
        Número certo                               & $50.6$ & $56.8$ & $+6.2$ \\ % [P]
        \hline
        Número errado                              & $12.5$ & $4.3$  & $-8.2$ \\ % [P]
        \hline
        Número atribuído sem número na anotação    & $2.5$  & $0.5$  & $-1.9$ \\ % [P]
        \hline
        Sem número, e a anotação também não tem    & $11.0$ & $12.9$ & $+1.9$ \\ % [P]
        \hline
        Sem número, mas a anotação tem             & $23.5$ & $25.5$ & $+2.0$ \\ % [P]
        \hline
    \end{tabular}
\end{table}

O efeito principal é a redução dos números errados para cerca de um terço, de $12.5\%$ para
$4.3\%$ das deteções, e dos números atribuídos a jogadores sem número de $2.5\%$ para $0.5\%$. % [P]
Ao mesmo tempo, o número certo sobe $6.2$ pontos percentuais. % [P]
Este resultado, que à partida parece contraditório com uma regra mais conservadora, explica-se pelo
filtro de qualidade mais permissivo: como a confiança passou a impedir as leituras inseguras, foi
possível admitir recortes mais pequenos e menos nítidos, e parte deles contém um número legível que
a primeira versão nem chegava a ver. O custo é o previsto no desenho da regra. Os jogadores que
ficam sem número apesar de o terem aumentam $2.0$ pontos, sobretudo casos em que o modelo leu o % [P]
número certo em várias janelas, mas sempre com confiança abaixo do limiar.

No GS-HOTA oficial, a subida é consistente nos três subconjuntos
(Tabela~\ref{tab:jnr-confianca-gshota}), incluindo o de treino, que não interveio em nenhuma decisão.

\begin{table}[htbp]
    \centering
    \caption{GS-HOTA oficial com a primeira versão do módulo e com a decisão pela confiança da
    leitura, sobre os mesmos vídeos ($117$). % [P]
    }
    \label{tab:jnr-confianca-gshota}
    \begin{tabular}{|l|c|c|c|}
        \hline
        \textbf{Conjunto} & \textbf{Primeira versão} & \textbf{Com confiança} & \textbf{$\Delta$} \\
        \hline
        Teste ($41$ vídeos)     & $43.01$ & $47.29$ & $+4.29$ \\ % [P]
        \hline
        Validação ($37$ vídeos) & $35.55$ & $37.19$ & $+1.64$ \\ % [P]
        \hline
        Treino ($39$ vídeos)    & $41.91$ & $45.00$ & $+3.09$ \\ % [P]
        \hline
    \end{tabular}
\end{table}

A melhoria não é uniforme. O GS-HOTA desce em $18$ dos $117$ vídeos, na maioria por menos de um % [P]
ponto. A inspeção dos casos com maior descida mostra dois padrões. No primeiro, o modelo lê o número
certo de forma repetida mas sempre com pouca confiança; no SNGS-056, um jogador com o número 14 foi
lido como 14 em nove de doze janelas, todas com confiança entre $0.3$ e $0.43$, e ficou sem número,
quando a primeira versão o acertava. É o custo direto do limiar. No segundo, a descida resulta da
própria anotação. No SNGS-188, um \emph{tracklet} lido como 14 com confiança entre $0.6$ e $0.8$ em
quase todas as janelas está emparelhado com um jogador anotado com o número 34, e os recortes
mostram de forma legível um 14. Casos semelhantes de anotação incompleta ou inconsistente foram
encontrados noutros vídeos, em que um número claramente visível não consta da anotação.

No extremo oposto, o SNGS-196 ilustra o comportamento pretendido. A primeira versão atribuía o
número certo a $9\%$ das deteções e um número errado a $16\%$; com a decisão pela confiança, o
número certo sobe para $50\%$ e o número errado desaparece, e o GS-HOTA do vídeo passa de $17.1$
para $38.7$. Dos quinze jogadores cujo número muda, doze que estavam sem número passam a ter o
número certo, dois números errados são corrigidos e um terceiro é retirado. Nenhum jogador que
estava correto passa a errado.


% -------------------------------------------------------------------------------------
% BLOCO E — Cap. 4. NOVA subsecção a seguir à anterior: alternativas testadas e rejeitadas.
% -------------------------------------------------------------------------------------

\subsection{Alternativas Testadas e Rejeitadas}\label{sec:res_jnr_rejeitadas}

Antes e depois de chegar à decisão pela confiança, foram avaliadas outras formas de reduzir o erro
do número. Todas foram medidas com o mesmo critério, o GS-HOTA oficial por subconjunto e vídeo a
vídeo, e só seriam aceites se não piorassem nenhum dos indicadores de erro. Nenhuma cumpriu esse
critério.

\paragraph{Fusão de fragmentos do seguimento.} Um mesmo jogador é frequentemente dividido pelo
\emph{tracker} em vários \emph{tracklets}, alguns dos quais demasiado curtos para serem lidos.
Juntar fragmentos da mesma equipa que nunca coexistem no tempo e são próximos no espaço de
reidentificação, propagando o número do fragmento mais confiante, aumentou a exatidão por
\emph{tracklet} em $5.5$ pontos percentuais. O GS-HOTA subiu apenas $0.04$ pontos. Os fragmentos
corrigidos eram quase todos curtos e representavam poucas deteções, e as fusões erradas
degradavam a associação entre identidades.

\paragraph{Regras de decisão alternativas.} Somar as confianças de cada número, em vez de contar
as janelas confiantes, não superou a regra final. Uma regra híbrida, que recorria à decisão da
primeira versão quando nenhuma janela era confiante, desde que houvesse pelo menos duas janelas
concordantes, pretendia recuperar os casos de leitura certa com pouca confiança. Numa amostra de
$20$ vídeos não trouxe ganho no GS-HOTA face à regra final e, no conjunto de teste, voltou a
atribuir números a jogadores sem número, que passaram de $0.1\%$ para $1.1\%$ das deteções.

\paragraph{Unicidade do número dentro da equipa.} Dois jogadores da mesma equipa não podem ter o
mesmo número em campo ao mesmo tempo. A inspeção visual mostrou casos em que esta regra é violada,
como no SNGS-196, em que um \emph{tracklet} perde o jogador que seguia, reaparece colado a outro
jogador da mesma equipa e leva consigo o número do primeiro. Testaram-se duas formas de impor a
regra: retirar o número ao \emph{tracklet} mais fraco nos \emph{frames} em que dois
\emph{tracklets} da mesma equipa partilham o número, e, numa versão mais restrita, fazê-lo apenas
quando os dois estão a mais de $5$~m de distância no campo, retirando o número ao que apareceu mais
recentemente. Numa primeira contagem em $60$ vídeos, dos $107$ conflitos encontrados, $68$ envolviam
dois \emph{tracklets} corretos, o mesmo jogador dividido em dois \emph{tracklets} sobrepostos por
alguns \emph{frames}. As duas versões deixaram o GS-HOTA igual ou mais baixo nos três subconjuntos, com descidas até % [P]
$0.15$ pontos, e pioraram mais de vinte vídeos. Mesmo no SNGS-196, retirar o número ao troço % [P]
errado baixou o GS-HOTA do vídeo, porque a métrica penaliza uma identidade cujo número muda a meio
mais do que ganha no troço corrigido. O erro de fundo é uma troca de identidade do seguimento, e
não se corrige do lado do número sem alterar o \emph{tracker}, que está fora do âmbito deste
trabalho.

\paragraph{Lição sobre a métrica.} Estas experiências mostram que ganhos medidos por
\emph{tracklet} não se traduzem necessariamente em ganhos no GS-HOTA, que pondera por deteção e
avalia a coerência das identidades ao longo do tempo. Uma correção que atue sobre \emph{tracklets}
curtos é praticamente invisível na métrica oficial, e uma correção que altere o número de parte de
um \emph{tracklet} pode prejudicá-la mesmo quando está certa. Por esta razão, todas as decisões
finais do módulo foram tomadas com a métrica oficial e não com indicadores intermédios.


% -------------------------------------------------------------------------------------
% BLOCO F — Cap. 4, "Teto da Abordagem". SUBSTITUI o último parágrafo ("O diagnóstico acabou
% por ser o mesmo..."), que deixou de ser verdade: a margem estava, em parte, na decisão.
% -------------------------------------------------------------------------------------

O diagnóstico foi o mesmo nas três experiências. Nos casos difíceis, o modelo falha a leitura
porque a informação não está presente nos \emph{pixels}, seja pela distância da câmara, seja pela
orientação do jogador, e nenhuma reponderação heurística, ajuste de contraste ou recombinação de
vistas cria informação onde ela não existe. O que estas experiências não resolviam era outro
problema: o sistema não sabia distinguir as leituras em que o modelo tinha visto o número das
leituras em que o tinha adivinhado. Essa informação não estava na imagem, mas no próprio modelo,
e a sua utilização na decisão, apresentada na Secção~\ref{sec:res_jnr_confianca}, foi o que
permitiu reduzir o erro sem alterar a leitura.


% -------------------------------------------------------------------------------------
% O QUE MUDA NO RESTO DA TESE QUANDO HOUVER OS 159 VÍDEOS (não é texto da tese):
%  - tab:jnr (exatidão por tracklet) e tab:jnr-distribuicao: recalcular com o sistema final;
%  - GS-HOTA por subconjunto (sec:res_gshota), decomposição por atributo, anexo vídeo a vídeo,
%    Resumo/Abstract, Introdução e Conclusão: números do sistema final;
%  - Custo computacional: a releitura com confiança acrescenta ~4 min por vídeo (medido: 82 min
%    para 20 vídeos); dizer que pode ser integrada na primeira leitura (a confiança sai da mesma
%    chamada ao modelo), o que eliminaria a segunda passagem;
%  - tab:jnr-categorias e tab:jnr-confianca-gshota: trocar [P] pelos valores dos 159;
%  - Conclusão/trabalho futuro: corte dos tracklets nas trocas de identidade (tracker) como a
%    origem dos conflitos de número; limiar de confiança por comprimento do tracklet.
% -------------------------------------------------------------------------------------
